% !TeX encoding = UTF-8
\documentclass{../plantilla/hvtbook}
\HVTTitulo{Hogares eficientes}
\HVTSubtitulo{Metodología para transformar una vivienda existente o diseñar una nueva a partir de las necesidades de sus habitantes, el clima y mediciones reales. Primero se protegen la salud, la seguridad y el confort; después se reduce la demanda de agua y energía, se incorporan tecnologías apropiadas y se verifican los resultados.}
\HVTNivel{Vivienda sostenible en Boyacá}
\HVTSerie{Publicación técnica}
\HVTNumero{05}
\HVTAccion{Explorar capítulos}
\HVTDescripcion{Proyecto de vivienda sostenible y hogares eficientes en Boyacá.}
\HVTCanonical{https://hidrogenoverdeturquesa.com/proyecto-hogares-eficientes}
\HVTRegreso{Volver a proyectos}{/\#portfolio}
\HVTPortada{../../images/portfolio/Eficiente.png}{Figura 1. Cada hogar requiere un diagnóstico propio de personas, clima, vivienda y consumos.}
\begin{document}
\HVTMakeTitle

\begin{HVTPage}{Contenido}{Ruta del proyecto}{Una vivienda que funciona para quienes la habitan}
\HVTLead{Un hogar eficiente no es una colección de dispositivos. Es una vivienda segura, saludable, asequible y adaptada a su contexto, capaz de prestar sus servicios con menor consumo y menor impacto.}
\begin{HVTTaskOrdered}{Capítulos}
\item Personas, usos y prioridades. \item Diagnóstico de la vivienda. \item Clima y relación con el lugar. \item Confort, salud y calidad del aire. \item Envolvente y estrategias pasivas. \item Inventario y gestión de la energía. \item Instalaciones eléctricas e iluminación. \item Agua, saneamiento y reúso seguro. \item Materiales, residuos y mantenimiento. \item Tecnologías complementarias. \item Plan de intervención y financiación. \item Medición, verificación y aprendizaje.
\end{HVTTaskOrdered}
\end{HVTPage}

\begin{HVTPage}{Capítulo 1}{Habitantes}{Personas, usos y prioridades}
\HVTLead{La eficiencia debe mejorar la vida cotidiana sin imponer soluciones que la familia no pueda usar, mantener o pagar. Las guías de la OMS relacionan la vivienda con temperatura interior, calidad del aire, seguridad, accesibilidad y salud. \HVTCite{1}}
\begin{HVTTable}{Pregunta}{Información}{Decisión}
\HVTTableRow{¿Quiénes habitan?}{Edad, movilidad, salud, permanencia y actividades}{Confort, accesibilidad y seguridad}
\HVTTableRow{¿Cómo se usa?}{Horarios, cocina, estudio, trabajo, cuidado y descanso}{Zonificación, iluminación y ventilación}
\HVTTableRow{¿Qué falla?}{Frío, calor, humedad, humo, ruido, goteras, costos y cortes}{Prioridad de intervención}
\HVTTableRow{¿Qué puede sostenerse?}{Presupuesto, conocimientos, repuestos y mantenimiento}{Tecnología apropiada}
\end{HVTTable}
\begin{HVTWarning}{Puerta de decisión 1}Primero se corrigen los riesgos para la vida, la salud y la habitabilidad; el ahorro económico no justifica mantener una condición insegura.\end{HVTWarning}
\end{HVTPage}

\begin{HVTPage}{Capítulo 2}{Diagnóstico}{Estado actual de la vivienda}
\HVTLead{La línea base combina inspección, planos, facturas, mediciones y conversación con los habitantes. Debe diferenciar problemas estructurales, constructivos, sanitarios, eléctricos y de uso.}
\begin{HVTTask}{Inspección inicial}\item Propiedad, uso permitido, área, antigüedad, reformas y planos disponibles. \item Estructura, cubierta, muros, pisos, puertas y ventanas. \item Humedad, filtraciones, moho, plagas, asbesto u otros materiales de preocupación. \item Instalaciones eléctricas, gas, abastecimiento de agua y saneamiento. \item Orientación, sombras, ventilación, iluminación natural y ruido. \item Facturas de energía y agua de al menos un periodo representativo.\end{HVTTask}
\HVTEquation{I_E=\frac{E_{\text{periodo}}}{A\cdot n_{\text{meses}}}}{}
\HVTText{La intensidad energética puede expresarse en kWh/m²·mes. También conviene calcular consumo por persona; ambos indicadores requieren contexto y no deben usarse para comparar hogares con servicios diferentes.}
\begin{HVTWarning}{Puerta de decisión 2}Si existen daños estructurales, riesgo eléctrico, combustión insegura, agua no apta o saneamiento deficiente, intervienen profesionales competentes antes de instalar equipos de eficiencia.\end{HVTWarning}
\end{HVTPage}

\begin{HVTPage}{Capítulo 3}{Clima}{Diseñar para el lugar}
\HVTLead{Colombia requiere estrategias diferenciadas para climas fríos, templados, cálidos secos y cálidos húmedos. El Atlas Climatológico del IDEAM y la clasificación climática de MinVivienda aportan una primera referencia, que debe complementarse con observación y medición local. \HVTCite{2} \HVTCite{3}}
\begin{HVTTable}{Variable}{Cómo observarla}{Aplicación}
\HVTTableRow{Temperatura y humedad}{Serie horaria interior y exterior}{Ventilación, sombra, aislamiento y control de humedad}
\HVTTableRow{Sol}{Recorrido, obstáculos y radiación disponible}{Orientación, aleros, iluminación y energía solar}
\HVTTableRow{Viento}{Dirección, velocidad y exposición}{Ventilación cruzada y protección}
\HVTTableRow{Lluvia}{Intensidad, estacionalidad y escorrentía}{Cubierta, drenaje y captación}
\HVTTableRow{Entorno}{Vegetación, construcciones, vías, humo y ruido}{Ubicación de aberturas y barreras}
\end{HVTTable}
\begin{HVTWarning}{Puerta de decisión 3}No se copia una solución de otro clima. Una abertura útil para enfriar puede aumentar pérdidas térmicas o introducir humedad, ruido y contaminación en otro lugar.\end{HVTWarning}
\end{HVTPage}

\begin{HVTPage}{Capítulo 4}{Salud y confort}{Aire limpio, temperatura y bienestar}
\HVTLead{La OMS recomienda abordar conjuntamente combustibles limpios, equipos de bajas emisiones, ventilación y diseño de la vivienda. Ventilar humo hacia afuera no reemplaza la eliminación o control de su fuente. \HVTCite{4}}
\begin{HVTTable}{Riesgo}{Diagnóstico}{Jerarquía de control}
\HVTTableRow{Combustión}{Combustible, equipo, evacuación y monóxido de carbono}{Eliminar o sustituir fuente; después extracción y ventilación}
\HVTTableRow{Humedad y moho}{Filtración, condensación, capilaridad o actividad interior}{Corregir causa y secar; no ocultar con pintura}
\HVTTableRow{Calor o frío}{Temperatura operativa, horarios y percepción}{Sombra, envolvente, ventilación y sistema activo si hace falta}
\HVTTableRow{Ruido}{Fuente, horario y ruta de transmisión}{Control en fuente, barrera y aislamiento}
\HVTTableRow{Iluminación}{Tarea visual, luz natural, deslumbramiento y noche}{Distribución, control solar y luminaria adecuada}
\end{HVTTable}
\begin{HVTTask}{Plan de medición}\item Ubicación y altura documentada de sensores. \item Temperatura y humedad en periodos representativos. \item CO2 como indicador complementario de ventilación, cuando aplique. \item CO y material particulado cuando existan fuentes de combustión. \item Registro de percepción de los habitantes junto con los datos.\end{HVTTask}
\end{HVTPage}

\begin{HVTPage}{Capítulo 5}{Diseño pasivo}{Reducir la demanda desde la envolvente}
\HVTLead{Antes de aumentar la capacidad de refrigeración, calefacción o generación, se revisan cubierta, sombra, muros, ventanas, infiltraciones y ventilación natural. La Resolución 0194 de 2025 organiza estrategias pasivas y activas según clima y uso de la edificación. \HVTCite{3}}
\HVTEquation{\dot{Q}_{\text{trans}}=U\cdot A\cdot(T_{\text{ext}}-T_{\text{int}})}{}
\HVTText{La ecuación estima transferencia térmica estacionaria por conducción. No incluye por sí sola radiación solar, infiltración, ventilación, humedad, masa térmica ni cargas internas.}
\begin{HVTFlow}\HVTFlowItem{1}{Proteger}{Resolver lluvia, humedad, infiltraciones de aire no deseadas y puentes térmicos críticos.}\HVTFlowItem{2}{Moderar}{Aplicar sombra, cubierta reflectiva o aislada, vegetación y ventanas apropiadas al clima.}\HVTFlowItem{3}{Ventilar}{Aprovechar viento y diferencia térmica solo cuando el aire exterior y la seguridad lo permitan.}\end{HVTFlow}
\begin{HVTWarning}{Puerta de decisión 5}Comparar alternativas mediante simulación o medición cuando la intervención afecte significativamente temperatura, humedad o ventilación. Más aislamiento no siempre equivale a mayor confort en todos los climas.\end{HVTWarning}
\end{HVTPage}

\begin{HVTPage}{Capítulo 6}{Energía}{Inventario, medición y gestión de cargas}
\HVTLead{Las potencias de placa sirven para una primera estimación; las cargas importantes se miden cuando sea posible y se contrastan con la factura.}
\HVTEquation{E_{\text{total}}=\sum_i(P_i\cdot t_i\cdot f_i)}{}
\HVTText{P es potencia, t tiempo de uso y f un factor que representa ciclos o fracción de operación. Las unidades deben ser coherentes para obtener kWh.}
\begin{HVTTable}{Carga}{Dato necesario}{Alternativa}
\HVTTableRow{Refrigeración}{kWh/día, estado de sellos, ubicación y temperatura}{Mantenimiento, ajuste o equipo eficiente}
\HVTTableRow{Calentamiento de agua}{Litros, temperatura, energía y horario}{Reducir demanda, aislar y evaluar tecnología}
\HVTTableRow{Cocción}{Combustible, eficiencia, emisiones y ventilación}{Equipo y fuente más limpia y segura}
\HVTTableRow{Iluminación}{Potencia, horas, iluminancia y deslumbramiento}{Luz natural, LED y control por zonas}
\HVTTableRow{Electrónica y espera}{Consumo activo y en espera}{Gestión de uso sin desconectar servicios críticos}
\end{HVTTable}
\begin{HVTTaskOrdered}{Orden de intervención}\item Eliminar consumos innecesarios sin afectar bienestar. \item Mantener y ajustar equipos existentes. \item Sustituir equipos al final de su vida útil usando información RETIQ cuando aplique. \HVTCite{5} \item Automatizar únicamente funciones útiles y mantenibles. \item Dimensionar generación renovable con la demanda ya optimizada.\end{HVTTaskOrdered}
\end{HVTPage}

\begin{HVTPage}{Capítulo 7}{Electricidad}{Seguridad eléctrica e iluminación}
\HVTLead{Toda ampliación, remodelación o intervención eléctrica debe estar dirigida, supervisada y ejecutada por personas técnica y legalmente competentes, conforme al RETIE vigente. Los sistemas de iluminación deben revisar requisitos del RETILAP. \HVTCite{6} \HVTCite{7}}
\begin{HVTTask}{Revisión por personal competente}\item Capacidad del servicio, tableros, circuitos y conductores. \item Puesta a tierra, protecciones contra sobrecorriente y choque eléctrico. \item Tomacorrientes, empalmes, extensiones permanentes y puntos calientes. \item Zonas húmedas, cocina, baño, exterior y equipos de alta potencia. \item Protección contra sobretensiones y riesgo por descargas atmosféricas cuando aplique. \item Planos, rotulado, pruebas y certificación exigible.\end{HVTTask}
\HVTEquation{P_{\text{demanda}}=\sum_i(P_{\text{conectada},i}\cdot f_{\text{demanda},i})}{}
\HVTText{Los factores de demanda y criterios de diseño no se eligen de manera arbitraria; deben aplicarse según reglamentación, tipo de carga y condiciones de la instalación.}
\begin{HVTWarning}{Seguridad}Enchufes inteligentes, baterías, inversores o paneles no corrigen una instalación defectuosa. Primero se resuelven las no conformidades eléctricas.\end{HVTWarning}
\end{HVTPage}

\begin{HVTPage}{Capítulo 8}{Agua}{Abastecimiento, uso y saneamiento}
\HVTLead{Se protege primero la calidad del agua para consumo humano y la continuidad del saneamiento. Después se reducen fugas y demandas, y solo entonces se estudian captación y reúso.}
\HVTEquation{V_{\text{uso}}=\sum_j(q_j\cdot t_j\cdot n_j)}{}
\HVTText{El volumen por uso se estima a partir del caudal q, duración t y frecuencia n, y se contrasta con medidor o factura.}
\begin{HVTTable}{Etapa}{Acción}{Verificación}
\HVTTableRow{Diagnóstico}{Medidor, presión, fugas, calidad y continuidad}{Balance de agua y prueba de fugas}
\HVTTableRow{Reducción}{Reparar, ajustar hábitos y seleccionar aparatos eficientes}{Caudal o descarga y consumo posterior}
\HVTTableRow{Captación de lluvia}{Evaluar precipitación, cubierta, primera descarga, tanque y tratamiento}{Calidad compatible con uso y rebose seguro}
\HVTTableRow{Aguas grises}{Separar en origen, tratar, almacenar y distribuir según diseño}{Control sanitario, señalización y ausencia de conexión cruzada}
\HVTTableRow{Saneamiento}{Conectar o tratar según contexto urbano o rural}{Permisos, operación, mantenimiento y disposición segura}
\end{HVTTable}
\HVTEquation{V_{\text{lluvia}}=P\cdot A\cdot C}{}
\HVTText{P es precipitación, A área efectiva y C coeficiente de escorrentía. El volumen útil depende además de pérdidas, calidad, demanda, estacionalidad y capacidad del tanque.}
\begin{HVTWarning}{Puerta de decisión 8}No conectar agua lluvia o gris a redes de consumo humano. El reúso requiere barreras sanitarias, identificación de tuberías y un plan de operación.\end{HVTWarning}
\end{HVTPage}

\begin{HVTPage}{Capítulo 9}{Circularidad}{Materiales, residuos y mantenimiento}
\HVTLead{La mejora más durable suele ser conservar y reparar lo que aún funciona. Cada material nuevo debe evaluarse por salud, procedencia, vida útil, mantenimiento, desmontaje y destino final.}
\begin{HVTTable}{Decisión}{Preguntas}{Evidencia}
\HVTTableRow{Conservar}{¿Es seguro, reparable y funcional?}{Inspección y mantenimiento}
\HVTTableRow{Reparar}{¿La intervención prolonga su vida sin crear riesgos?}{Diagnóstico y garantía}
\HVTTableRow{Sustituir}{¿El beneficio supera impactos y costo del reemplazo?}{Comparación de ciclo de vida y desempeño}
\HVTTableRow{Desmontar}{¿Puede recuperarse sin contaminar otras corrientes?}{Plan de separación y almacenamiento}
\HVTTableRow{Disponer}{¿Es ordinario, aprovechable, especial o peligroso?}{Gestor y comprobante cuando aplique}
\end{HVTTable}
\begin{HVTTask}{Plan de mantenimiento}\item Cubierta, canales, bajantes y drenajes. \item Sellos de puertas y ventanas. \item Filtros, extractores y equipos de combustión. \item Tanques, bombas y sistemas de tratamiento. \item Tableros y protecciones eléctricas. \item Registro de fallas, repuestos y costos.\end{HVTTask}
\end{HVTPage}

\begin{HVTPage}{Capítulo 10}{Tecnologías}{Complementos después de reducir la demanda}
\HVTLead{La selección tecnológica depende del recurso local, el servicio requerido, la capacidad de la vivienda y el mantenimiento disponible.}
\begin{HVTTable}{Tecnología}{Estudio previo}{Riesgo que debe controlarse}
\HVTTableRow{Solar fotovoltaica}{Demanda horaria, recurso, sombras, cubierta y red}{Estructura, electricidad, incendio y trabajo en altura}
\HVTTableRow{Solar térmica}{Demanda de agua caliente, clima y calidad de agua}{Quemaduras, presión, estancamiento y legionela según sistema}
\HVTTableRow{Baterías}{Cargas críticas, autonomía y estrategia de operación}{Compatibilidad, ventilación, protección e incendio}
\HVTTableRow{Automatización}{Problema concreto, ahorro medible y uso manual}{Privacidad, ciberseguridad, dependencia y obsolescencia}
\HVTTableRow{Cocción limpia}{Uso, costo, confiabilidad y calidad del aire}{Combustión, electricidad y ventilación}
\HVTTableRow{Hidrógeno}{Necesidad energética, suministro, equipos y normativa}{No se incorpora en vivienda sin ingeniería especializada y gestión integral de seguridad}
\end{HVTTable}
\HVTEquation{E_{FV}\approx P_0\cdot H_{POA}\cdot PR}{}
\HVTText{Esta aproximación relaciona potencia nominal, irradiación sobre el plano y factor global de desempeño. El diseño final requiere serie meteorológica, sombras, temperatura, pérdidas, estructura, protecciones y normativa.}
\begin{HVTWarning}{Puerta de decisión 10}Cada tecnología debe demostrar utilidad, seguridad, repuestos, mantenimiento y costo total. “Inteligente” o “renovable” no significa automáticamente eficiente.\end{HVTWarning}
\end{HVTPage}

\begin{HVTPage}{Capítulo 11}{Implementación}{Plan de intervención y financiación}
\HVTLead{Las medidas se agrupan para evitar obras repetidas y proteger el presupuesto familiar. La estimación económica incluye diseño, instalación, permisos, mantenimiento, reposición y disposición final.}
\begin{HVTTable}{Fase}{Contenido}{Criterio de cierre}
\HVTTableRow{0. Emergente}{Riesgos estructurales, eléctricos, sanitarios o de combustión}{Condición segura documentada}
\HVTTableRow{1. Bajo costo}{Fugas, ajustes, mantenimiento, formación y medición}{Resultado verificado}
\HVTTableRow{2. Envolvente y redes}{Cubierta, humedad, sombra, ventanas e instalaciones}{Pruebas y recepción}
\HVTTableRow{3. Equipos}{Sustitución justificada y controles}{Desempeño y garantía}
\HVTTableRow{4. Renovables}{Generación y almacenamiento dimensionados}{Puesta en servicio segura}
\end{HVTTable}
\HVTEquation{PB_{\text{simple}}=\frac{I_{\text{inicial}}}{A_{\text{anual}}}}{}
\HVTText{El periodo simple de recuperación es solo un filtro preliminar. La decisión debe considerar vida útil, reposiciones, financiación, variación tarifaria, beneficios de salud y confort, y riesgos.}
\begin{HVTTask}{Expediente de obra}\item Alcance, planos, especificaciones y presupuesto. \item Responsables competentes, permisos y seguros. \item Cronograma compatible con la vida familiar. \item Protección de habitantes y bienes durante la obra. \item Pruebas, garantías, manuales y capacitación.\end{HVTTask}
\end{HVTPage}

\begin{HVTPage}{Capítulo 12}{Verificación}{Comprobar, ajustar y aprender}
\HVTLead{El proyecto no termina al instalar equipos. Se compara la situación posterior con una línea base equivalente, corrigiendo cambios de clima, ocupación, tarifas y nivel de servicio.}
\HVTEquation{\text{Ahorro (\%)}=\frac{C_{\text{base ajustada}}-C_{\text{posterior}}}{C_{\text{base ajustada}}}\cdot100}{}
\begin{HVTTable}{Resultado}{Indicador}{Frecuencia}
\HVTTableRow{Energía}{kWh/mes, kWh/persona y kWh/m²}{Mensual y anual}
\HVTTableRow{Agua}{m³/mes y L/persona·día}{Mensual}
\HVTTableRow{Confort}{Temperatura, humedad y percepción}{Por temporada}
\HVTTableRow{Aire interior}{Indicadores definidos según las fuentes}{Durante usos críticos}
\HVTTableRow{Economía}{Costo real, ahorro y mantenimiento}{Trimestral o anual}
\HVTTableRow{Confiabilidad}{Fallas, interrupciones y tiempo de reparación}{Por evento}
\end{HVTTable}
\end{HVTPage}

\begin{HVTPage}{Capítulo 12}{Verificación}{Entregable y cierre del ciclo}
\begin{HVTTask}{Entregable integral}\item Diagnóstico y línea base. \item Mapa climático, de confort y riesgos. \item Paquete priorizado de medidas. \item Diseños y memorias de cálculo aplicables. \item Presupuesto, fases y responsables. \item Manual de uso y mantenimiento. \item Informe de medición y verificación.\end{HVTTask}
\begin{HVTWarning}{Cierre del ciclo}Si una medida ahorra pero reduce confort, salud o confiabilidad, se revisa. El objetivo es prestar mejores servicios con menos recursos, no trasladar el costo a los habitantes.\end{HVTWarning}
\end{HVTPage}

\begin{HVTPage}{Fuentes}{Base técnica}{Referencias consultadas}
\begin{HVTReferences}
\HVTReference{Organización Mundial de la Salud. \emph{WHO Housing and health guidelines}}{https://www.who.int/publications/i/item/9789241550376}
\HVTReference{IDEAM. \emph{Atlas Climatológico de Colombia}}{https://www.ideam.gov.co/AtlasWeb/}
\HVTReference{Ministerio de Vivienda, Ciudad y Territorio. \emph{Resolución 0194 de 2025 y Guía de construcción sostenible}}{https://www.minvivienda.gov.co/normativa/resolucion-0194-2025}
\HVTReference{Organización Mundial de la Salud. \emph{Guidelines for indoor air quality: household fuel combustion}}{https://www.who.int/publications/i/item/9789241548885}
\end{HVTReferences}
\end{HVTPage}
\begin{HVTPage}{Fuentes}{Base técnica}{Referencias consultadas}
\begin{HVTReferences}
\HVTReference{UPME. \emph{Reglamento Técnico de Etiquetado con fines de Uso Racional de la Energía — RETIQ}}{https://bdigital.upme.gov.co/handle/001/1022}
\HVTReference{Ministerio de Minas y Energía. \emph{Reglamento Técnico de Instalaciones Eléctricas — RETIE}}{https://www.minenergia.gov.co/es/misional/energia-electrica-2/reglamentos-tecnicos/reglamento-t\%C3\%A9cnico-de-instalaciones-el\%C3\%A9ctricas-retie/}
\HVTReference{Ministerio de Minas y Energía. \emph{Reglamento Técnico de Iluminación y Alumbrado Público — RETILAP}}{https://www.minenergia.gov.co/es/misional/energia-electrica-2/reglamentos-tecnicos/reglamento-t\%C3\%A9cnico-de-iluminaci\%C3\%B3n-y-alumbrado-p\%C3\%BAblico-retilap/}
\end{HVTReferences}
\HVTText{Las normas pueden actualizarse. Cada intervención debe comprobar las versiones vigentes, el alcance aplicable, los requisitos municipales y la competencia de quienes diseñan, construyen e inspeccionan. Esta publicación contiene únicamente información autorizada para divulgación pública. Los resultados detallados, datos operativos, análisis económicos y demás información estratégica del proyecto son confidenciales.}
\end{HVTPage}

\begin{HVTPage}{Colaboración}{Conversemos}{¿Necesitas investigar, formular o evaluar un proyecto relacionado?}
\HVTLead{Conversemos sobre tu necesidad de diagnóstico, diseño o evaluación de una vivienda eficiente, saludable y sostenible.}
\HVTContact{Consultar proyecto de vivienda por WhatsApp}{https://wa.me/573209574884}
\end{HVTPage}
\end{document}
